\pdfvariable suppressoptionalinfo 611
\providecommand{\CRevisionRound}{2}
\documentclass[10pt]{article}
\usepackage[a4paper,margin=19mm]{geometry}
\usepackage{amsmath,amssymb,amsthm,booktabs,microtype}
\usepackage[T1]{fontenc}
\usepackage{lmodern}
\usepackage[hidelinks,hypertexnames=false]{hyperref}
\newtheorem{theorem}{Theorem}
\newtheorem{proposition}{Proposition}
\newcommand{\R}{\mathbb R}
\title{Uniform Mutation in a Three-Strategy Rock--Paper--Scissors Flow}
\author{HCS Research Program}
\date{29 August 2026\quad (revision \CRevisionRound)}
\begin{document}
\maketitle
\begin{abstract}
We give a convention-complete theorem for the cyclic three-strategy flow
\[
 \dot x=ax(y-z)+\mu(1/3-x),\qquad\text{cyclically in }(x,y,z),
 \qquad a,\mu\geq0,
\]
on the unit simplex.  At zero mutation and positive cyclic rate, the product
$H=xyz$ is conserved: every regular interior level is one periodic orbit, the
boundary is a cyclic heteroclinic network, and the physical period has an
exact turning-point quadrature with center limit $2\pi\sqrt3/a$.  Positive additive uniform
mutation points inward, makes $\log H$ a strict Lyapunov function, and gives
global convergence to the barycenter; the tangent rates are
$-\mu\pm ia/\sqrt3$.  We close the exactly solvable $a=0$ contraction and
identity faces and keep finite numerical rows separate from the theorem.
This is a source-local population-flow result: there is no arithmetic owner,
target determinant, or Hilbert--Pólya claim.
\end{abstract}

\section{Model and simplex geometry}
Let
\[
 \Delta=\{(x,y,z)\in\R^3:x,y,z\ge0,\ x+y+z=1\}.
\]
The vector field is
\begin{equation}
 F_{a,\mu}(x,y,z)=\bigl(ax(y-z)+\mu(1/3-x),
 ay(z-x)+\mu(1/3-y),
 az(x-y)+\mu(1/3-z)\bigr). \tag{1}
\end{equation}
The cyclic terms cancel, hence $\dot x+\dot y+\dot z=0$.  On a coordinate
face and $\mu=0$ the face is invariant.  If $\mu>0$, a zero coordinate has
derivative $\mu/3$, so every trajectory enters the open simplex immediately.
The barycenter $b=(1/3,1/3,1/3)$ is fixed for all parameters.

\begin{theorem}[Conservative product and periodic levels]
For $\mu=0$, $H=xyz$ is constant.  If $a>0$, the barycenter is the only
interior equilibrium, every level $0<h<1/27$ is a single regular periodic
orbit, and as $h\downarrow0$ its closure tends to the three-edge
heteroclinic network.  If $a=0$, the flow is the identity instead.
\end{theorem}
\begin{proof}
On the interior,
\[
 \frac{d}{dt}\log H=a[(y-z)+(z-x)+(x-y)]=0. \tag{2}
\]
The equations $y+z=1-x$ and $yz=h/x$ have two distinct roots except at the
barycenter.  Consequently the compact positive level is a one-dimensional
regular component.  For $a>0$, the cyclic orientation of (1) gives a
nonvanishing vector field on it, hence one periodic orbit.  Setting one
coordinate to zero gives the invariant edges and their cyclic connections;
these are the $h=0$ boundary limit.  For $a=0$ every point is fixed, as also
follows directly from (1).
\end{proof}

\section{Exact physical period}
Assume $a>0$.  From the two-root relation,
\[
 (y-z)^2=(1-x)^2-\frac{4h}{x}.
\]
Let $x_-<1/3<x_+$ be the physical roots of
\begin{equation}
 x(1-x)^2=4h,\qquad 0<h<1/27. \tag{3}
\end{equation}
The two signs of $y-z$ traverse the closed level, so its exact period is
\begin{equation}
 T(h)=\frac2a\int_{x_-}^{x_+}
 \frac{dx}{x\sqrt{(1-x)^2-4h/x}}. \tag{4}
\end{equation}
The cubic has a third root $x_3=2-x_--x_+>1$.  With
\[
 x(\theta)=\frac{x_-+x_+}{2}+\frac{x_+-x_-}{2}\sin\theta,
 \qquad -\frac\pi2\le\theta\le\frac\pi2,
\]
the endpoint cosine cancels exactly and (4) becomes
\begin{equation}
 T(h)=\frac2a\int_{-\pi/2}^{\pi/2}
 \frac{d\theta}{\sqrt{x(\theta)\,[x_3-x(\theta)]}}. \tag{5}
\end{equation}
This regular integral is the deterministic representation used by the ledger.

\begin{proposition}[Center and boundary limits]
For $a>0$, $T(h)\to2\pi\sqrt3/a$ as $h\uparrow1/27$, while
$T(h)\to+\infty$ as $h\downarrow0$.
\end{proposition}
\begin{proof}
At the center, both physical roots approach $1/3$ and $x_3$ approaches
$4/3$; dominated convergence in (5) gives the displayed constant.  At the
lower endpoint, the factor $x(\theta)^{-1/2}$ develops a divergent integral,
which is the familiar slowing near the heteroclinic edges.
\end{proof}

\section{Uniform mutation and convergence}
\begin{theorem}[Strict product Lyapunov law]
For $\mu>0$ and every interior point,
\begin{equation}
 \frac{d}{dt}\log(xyz)=\frac\mu3\left(\frac1x+\frac1y+\frac1z-9\right)\ge0, \tag{6}
\end{equation}
with equality exactly at $b$.  Every solution from $\Delta$ converges to $b$;
there is no nonconstant recurrent trajectory.
\end{theorem}
\begin{proof}
The reciprocal AM--HM inequality and $x+y+z=1$ give
$1/x+1/y+1/z\ge9$, with equality only when all coordinates are equal.  A
boundary initial condition enters the interior for positive time, after which
$H>0$ and (6) applies.  The compact positively invariant simplex supplies a
nonempty omega-limit set.  LaSalle's invariance principle puts it in the
zero-derivative set, which is the singleton $\{b\}$.
\end{proof}

At $b$, the ambient Jacobian has characteristic polynomial
\[
 (\lambda+\mu)\left((\lambda+\mu)^2+\frac{a^2}{3}\right).
\]
Thus the tangent-plane eigenvalues are $-\mu\pm ia/\sqrt3$.  The imaginary
part is the residual cyclic rotation; the negative real part is supplied only
by the uniform mutation term.

\ifnum\CRevisionRound>0
\section{Degenerate faces and finite receipt}
When $a=0$, the three coordinates decouple:
\begin{equation}
 x_i(t)=\frac13+\bigl(x_i(0)-\tfrac13\bigr)e^{-\mu t}. \tag{7}
\end{equation}
When $a=\mu=0$, (1) is the identity on all of $\Delta$.  For $\mu=0,a>0$,
the edges remain invariant and the interior period family above is the whole
recurrence statement.  We make no extension to a general mutation matrix.

The receipt fixes three positive conservative rates, five rational product levels, three
center-limit rows, six mutation starts (including all three coordinate
boundaries), three exact contractions, and four tangent spectra.  Fixed-step
RK4 rows are explicitly labelled diagnostics.  The independent checker does
not import the producer; SymPy checks sixteen identities, byte replay uses two
fresh temporary trees, and the hostile suite rejects 25/25 mutations,
including a repaired-hash attempt to erase the positive-rate condition.
\fi

\ifnum\CRevisionRound>1
\section{Route-A boundary and audit supplement}
The flow has no intrinsic rational-prime carrier, no prime-power repetition
clock, no weighted primitive-orbit product, and no target divisor.  Its finite
periods are physical ODE periods on a continuum and are not an Artin--Mazur
zeta.  Accordingly the strict tuple is
\[
(\mathtt{A0\_FAIL},\mathtt{A1\_WEAK},\mathtt{A2\_FAIL},
 \mathtt{A3\_FAIL},\mathtt{A4\_FORMAL\_HINT}),
\]
with \texttt{ROUTE\_A\_REJECTED} and
\texttt{route\_b\_invocation\_allowed=false}.

The algebraic AM--HM remainder used by the checker is
\[
\frac1x+\frac1y+\frac1z-\frac9{x+y+z}
 =\frac{(x-y)^2z+(y-z)^2x+(z-x)^2y}{xyz(x+y+z)}\ge0.
\]
The finite rows certify implementation identities and mass closure; they do
not replace the compactness/LaSalle proof or establish any external spectral
matching.  The scope literal is
\texttt{NO\_BAD\_EULER\_OR\_ROOT\_NUMBER}.
\fi

\begin{thebibliography}{9}
\bibitem{mayleonard} R. M. May and W. J. Leonard, Nonlinear aspects of
competition between three species, \emph{SIAM Journal on Applied Mathematics}
29, 243--253 (1975).  This citation supplies cyclic-competition context only;
the frozen model here is the additive-uniform-mutation replicator field.
\bibitem{hofbauer} J. Hofbauer and K. Sigmund, \emph{Evolutionary Games and
Population Dynamics}, Cambridge University Press (1998).
\bibitem{lasalle} J. P. LaSalle, Some extensions of Liapunov's second method,
\emph{IRE Transactions on Circuit Theory} 7 (1960).
\end{thebibliography}

\section*{Declarations}
\textbf{Scope.} \texttt{NO\_BAD\_EULER\_OR\_ROOT\_NUMBER}; no target
arithmetic, target-zero, determinant-matching, or Hilbert--Pólya claim.
\textbf{Data and code.} All rows and audits accompany HCS-C235.  This is not
external peer review.
\textbf{AI-use disclosure.} Generative tools assisted drafting and code
generation; the internal artifact chain checks displayed formulas and metadata.
\end{document}
